\chapter{Implementierung des Proof of Concept} \label{kap:6}

Dieses Kapitel beschreibt die prototypische Umsetzung des in Kapitel~\ref{kap:5} konzipierten Assistenzsystems. Zunächst werden der technische Rahmen und die Projektstruktur dargestellt. Anschließend wird die Implementierung der einzelnen Pipeline-Module erläutert. Abschließend wird die Umsetzung der Prompt-Strategie einschließlich des System-Prompts beschrieben.

% ─────────────────────────────────────────────────────────────────────────────
\section{Technischer Rahmen und Projektstruktur}
\label{sec:techstack}

Der \ac{PoC} ist als Node.js-Anwendung auf Basis von TypeScript implementiert.\footcite[Vgl.][]{JavaScriptSyntaxTypesa} Die Wahl von TypeScript ist durch die technologische Nähe zum Untersuchungsgegenstand begründet: Der \ac{BWEC} ist in ein TypeScript-basiertes Monorepo eingebettet, sodass auch der \ac{PoC} im selben technologischen Kontext entwickelt wird. Dies erleichtert sowohl die Integration als auch eine mögliche Übertragung auf weitere Anwendungen innerhalb vergleichbarer Architekturen. Darüber hinaus ermöglicht TypeScript eine typsichere Definition des \ac{UFS} (vgl. Listing~\ref{lst:ufs}). Als Laufzeitumgebung wird Node.js ab Version~24 eingesetzt; die Ausführung erfolgt über \texttt{tsx}, einen TypeScript-Runner, der Kompilierung und Ausführung in einem Schritt verbindet. Auf ein zusätzliches Framework wurde bewusst verzichtet, um den \ac{PoC} minimal zu halten.

Als Kernabhängigkeiten dienen \texttt{playwright}@1.50 für die Browser-Automatisierung,\footcite[Vgl.][]{playwright} \texttt{@axe-core/playwright}@4.10 für die axe-core-Integration sowie \texttt{dotenv}@17.3 für die Konfigurationsverwaltung und \texttt{tsx}@6.5.7 als TypeScript-Runner.

Die Projektstruktur folgt einer flachen Modulorganisation (vgl.\,Anhang~\ref{app:projektstruktur}). Jedes Modul in \texttt{src/modules/} übernimmt sowohl die Datenerhebung als auch die Normalisierung in das \ac{UFS}.

Die Pipeline wird über \texttt{src/index.ts} orchestriert und unterstützt fünf \ac{CLI}-Flags: \texttt{-{}-url} (Ziel-\acs{URL}), \texttt{-{}-src-dir} (Quellcode-Pfad), \texttt{-{}-skip-llm} (Prompt speichern ohne finale Priorisierung), \texttt{-{}-axe-only} (nur axe-core ausführen) und \texttt{-{}-llm-detect} (zusätzliche \ac{LLM}-Codeanalyse aktivieren). Die Konfiguration -- einschließlich Ollama-\acs{URL}, Modellname und Timeout-Werte -- erfolgt zentral in \texttt{config.ts} und kann über Umgebungsvariablen überschrieben werden.

% ─────────────────────────────────────────────────────────────────────────────
\section{axe-core-Modul}
\label{sec:impl-axe}

Das axe-core-Modul realisiert \ac{FA}-01 und übernimmt die regelbasierte Accessibility-Erkennung. Es startet über Playwright einen headless Chromium-Browser, navigiert zur Ziel-\acs{URL} und wartet auf \texttt{networkidle}, um sicherzustellen, dass die React-Anwendung vollständig gerendert ist. Anschließend wird über \texttt{AxeBuilder} eine axe-core-Analyse ausgeführt, die auf die Tags \texttt{wcag2a}, \texttt{wcag2aa}, \texttt{wcag21a}, \texttt{wcag21aa}, \texttt{wcag21a} und \texttt{wcag22aa} beschränkt ist -- entsprechend dem gesetzlichen Mindeststandard (vgl. Abschnitt~\ref{kap:2.1}).

Die Normalisierung überführt jede Violation in ein \texttt{UnifiedFinding}. Dabei erfolgen drei zentrale Transformationen:

\begin{itemize}
  \item \textbf{Severity-Mapping:} Das \texttt{impact}-Feld von axe-core (\texttt{critical}, \texttt{serious}, \texttt{moderate}, \texttt{minor}) wird direkt auf das \texttt{severity}-Feld des \ac{UFS} abgebildet.
  \item \textbf{\ac{WCAG}-Extraktion:} Aus den axe-Tags (z.\,B. \texttt{wcag143}) werden die Erfolgskriterien (z.\,B. \texttt{1.4.3}) und die Konformitätsstufe extrahiert.
  \item \textbf{Kategorie-Heuristik:} Jede Regel wird anhand ihres \texttt{ruleId} einer Kategorie der in Abschnitt~\ref{Taxonomie} eingeführten Taxonomie zugeordnet. Farbkontrast-Regeln werden als \texttt{layout} klassifiziert, ARIA-Strukturregeln als \texttt{semantisch}, alle übrigen als \texttt{syntaktisch}.
\end{itemize}

Die Felder \texttt{componentPath}, \texttt{codeSnippet} und \texttt{lineRange} bleiben nach der axe-Normalisierung zunächst \texttt{null} und werden erst in der Code-Phase (Abschnitt~\ref{sec:impl-code}) angereichert.

Zu berücksichtigen ist, dass axe-core im \ac{PoC} nur einmal beim initialen Seitenaufruf ausgeführt wird und nicht nach jeder Playwright-Interaktion. Eine Analyse pro Zustandswechsel wäre technisch realisierbar, würde jedoch bei sieben Checks mit je mehreren Zustandsänderungen die Laufzeit um bis zu 168~Sekunden erhöhen und damit \ac{NFA}-02 verletzen. Diese Einschränkung wird in Abschnitt~\ref{sec:limitationen} als Limitation diskutiert.

% ─────────────────────────────────────────────────────────────────────────────
\section{Playwright-Modul}
\label{sec:impl-pw}

Das Playwright-Modul implementiert \ac{FA}-02 (dynamische Zustandserfassung). Es definiert sieben generische Interaktionschecks, die jeweils ein spezifisches \ac{WCAG}-Erfolgskriterium adressieren (vgl. Anhang~\ref{app:pw_grep_checks}, Tabelle~\ref{tab:pw-checks}).

Die Checks sind bewusst anwendungsunabhängig konzipiert: Sie setzen keine Kenntnis der konkreten Anwendungslogik voraus und sind auf jeder Web-App ausführbar. Der Keyboard-Trap-Check beispielsweise simuliert 40 Tab-Schritte und prüft, ob der Fokus in den letzten zehn Schritten auf demselben Element verbleibt. Der Check \texttt{focus-visible} navigiert per Tab durch die Seite und prüft für jedes fokussierte Element, ob ein sichtbarer Fokusindikator (\texttt{outline} oder \texttt{box-shadow}) vorhanden ist -- eine Prüfung, die Tafreshipour et al. als Schwachstelle bestehender Tools identifizieren.\footcite[Vgl.][S.~110]{tafreshipour_ma11y_2024}

Fehlgeschlagene Checks werden über ein statisches Mapping (\texttt{CHECK\_DEFINITIONS}) in \texttt{UnifiedFinding}-Objekte überführt. Jede Check-Definition enthält die zugehörige \texttt{ruleId}, \texttt{severity}, \ac{WCAG}-Kriterien und Kategorie. Checks, die erfolgreich bestanden werden, erzeugen kein Finding.

% ─────────────────────────────────────────────────────────────────────────────
\section{Quellcode-Anreicherung und Anti-Pattern-Erkennung}
\label{sec:impl-code}

Das Code-Modul übernimmt zwei Aufgaben: die Anreicherung bestehender Findings mit Quellcode-Kontext (\ac{FA}-03) und die eigenständige Erkennung von Anti-Patterns im React-Quellcode.

\subsection{Anreicherung bestehender Findings}

Für jedes axe-core- oder Playwright-Finding, das einen CSS-Selektor enthält, aber noch keinen \texttt{componentPath} besitzt, wird der Quellcode der React-Anwendung durchsucht. Der Algorithmus extrahiert Suchbegriffe aus dem Selektor -- IDs, CSS-Klassen, ARIA-Attribute -- und sucht diese in den \texttt{.tsx}/\texttt{.jsx}/\texttt{.ts}/\texttt{.js}-Dateien des angegebenen Quellverzeichnisses. Bei einem Treffer werden \texttt{componentPath}, \texttt{codeSnippet} (zehn Zeilen Kontext um die Fundstelle) und \texttt{lineRange} im Finding ergänzt.

Diese Anreicherung adressiert die in Abschnitt~\ref{Herausforderungen} beschriebene Diskrepanz zwischen gerendertem \ac{DOM} und Quellcode. Durchsucht werden ausschließlich Dateien, die durch axe-core- oder Playwright-Findings referenziert werden -- nicht die gesamte Codebasis. Dieses selektive Vorgehen reduziert den Token-Verbrauch im späteren Prompt (vgl. Abschnitt~\ref{sec:kontextualisierung}).

\subsection{Pattern-Findings}

Zusätzlich zur Anreicherung erkennt das Modul sechs Anti-Patterns durch reguläre Ausdrücke (vgl. Anhang~\ref{app:pw_grep_checks}, Tabelle~\ref{tab:grep-patterns}): \texttt{div-span-onclick}, \texttt{img-missing-alt}, \texttt{label-missing-htmlfor}, \texttt{role-button-non-semantic}, \texttt{tabindex-removes-element} und \texttt{autofocus-usage}.

Treffer werden nach der Normalisierung dedupliziert und auf maximal acht Treffer pro Pattern begrenzt, um eine Überrepräsentation im Token-Budget zu vermeiden (vgl. Abschnitt~\ref{sec:kontextualisierung}).

% ─────────────────────────────────────────────────────────────────────────────
\section{\acs{LLM}-Detektor-Modul}
\label{sec:impl-llm-detector}

Das \acs{LLM}-Detektor-Modul ergänzt die regelbasierten Quellen um eine modellbasierte Quellcodeanalyse. Es wird nur ausgeführt, wenn der \ac{CLI}-Parameter \texttt{-{}-llm-detect} gesetzt ist und ein Quellcodepfad via \texttt{-{}-src-dir} vorliegt. Ziel ist nicht die Ersetzung bestehender Tools, sondern die Ergänzung um eine zusätzliche Detektionsquelle mit semantischem Fokus.

Die Verarbeitung erfolgt in fünf Schritten:

\begin{enumerate}
  \item \textbf{Dateiselektion:} Das Modul sammelt rekursiv Dateien mit den Endungen \texttt{.tsx}, \texttt{.ts}, \texttt{.jsx}, \texttt{.js} und \texttt{.css}. Die Anzahl der analysierten Dateien ist über \texttt{LLM\_DETECT\_MAX\_FILES} begrenzbar (Default: 60; in den Benchmarkläufen auf 25 reduziert).
  \item \textbf{Chunking:} Die Dateien werden mit Zeilennummern versehen und in Text-Chunks aufgeteilt. Die Chunk-Größe ist über \texttt{LLM\_DETECT\_CHUNK\_CHARS} konfigurierbar (Benchmarkläufe: 4\,000 Zeichen).
  \item \textbf{Detektionsprompt:} Für jeden Chunk wird ein eigener Prompt erzeugt, der ausschließlich \acs{JSON} im fest definierten Schema verlangt.\footnote{Der vollständige System- und User-Prompt der Detektor-Phase ist in Anhang~\ref{app:detektorprompt} dokumentiert.}
  \item \textbf{Parsing und Validierung:} Antworten werden als \acs{JSON} extrahiert und validiert; ein Fallback überführt Formatabweichungen in das Zielformat.
  \item \textbf{Normalisierung und Deduplizierung:} Valide Einträge werden als \texttt{source="llm"} in das \ac{UFS} überführt und über den Schlüssel \texttt{ruleId|filePath|startLine} dedupliziert.
\end{enumerate}

Dieses Vorgehen reduziert das Risiko von Halluzinationen, weil Findings ohne belastbare Evidenz oder ohne verwertbaren Datei-/Zeilenbezug verworfen werden. Gleichzeitig bleibt die Nachvollziehbarkeit erhalten, da jedes \ac{LLM}-Finding einem konkreten Quellcodeausschnitt zugeordnet ist.

\subsection{Abgrenzung zu grep und Nutzen der Kombination}

Die beiden Detektionsansätze unterscheiden sich grundlegend in ihrer Arbeitsweise:

\begin{itemize}
  \item \textbf{grep} arbeitet regelbasiert und deterministisch. Es erkennt bekannte Muster zuverlässig und reproduzierbar, kann aber nur das finden, was vorher als Regel modelliert wird.
  \item \textbf{\acs{LLM}-Detektor} arbeitet kontextsensitiv. Er kann auch semantisch ähnliche Problemstellen erkennen, die nicht exakt einem festen Regex-Muster entsprechen, ist jedoch anfälliger für Formatabweichungen und Halluzinationen.
\end{itemize}

Der \ac{PoC} kombiniert daher beide Ansätze in komplementärer Weise: grep als stabile Baseline und den \acs{LLM}-Detektor als zusätzliche Quelle mit potenziellem Recall-Gewinn.

\subsection{Beispielhafter Ablauf eines \acs{LLM}-Findings}

Ein typischer Detektionsfall läuft wie folgt ab: In einer Komponente enthält ein \texttt{<div>} einen \texttt{onClick}-Handler ohne tastaturäquivalentes Ereignis. Das Modell meldet dazu ein Finding mit Datei- und Zeilenbezug sowie Evidenztext, das normalisiert, dedupliziert und als \texttt{source="llm"} in das \ac{UFS} überführt wird, bevor es in der Priorisierungsphase verarbeitet wird.

\begin{lstlisting}[language={}, caption={Beispiel eines validierten \acs{LLM}-Detektor-Findings (vereinfacht)}]
{
  "title": "Klickbare div-Komponente ohne Tastaturbedienung",
  "ruleId": "keyboard-click-only",
  "wcagCriteria": ["2.1.1"],
  "wcagLevel": "A",
  "severity": "serious",
  "category": "semantisch",
  "filePath": "components/ContactForm.tsx",
  "startLine": 45,
  "endLine": 49,
  "evidence": "<div ... onClick={() => setTopic(t)}>",
  "fixSuggestion": "button verwenden oder onKeyDown fuer Enter/Space ergaenzen"
}
\end{lstlisting}

% ─────────────────────────────────────────────────────────────────────────────
\section{Prompt-Builder und System-Prompt}
\label{sec:impl-prompt}

Der Prompt-Builder setzt die in Abschnitt~\ref{sec:kontextualisierung} beschriebene Prompt-Strategie technisch um. Er besteht aus drei Komponenten: dem System-Prompt, der Token-Budget-Steuerung und der Serialisierung der Findings aus vier Quellen (axe-core, Playwright, grep, \ac{LLM}-Codeanalyse).

\subsection{System-Prompt}

Der System-Prompt implementiert die in Kapitel~\ref{kap:5} gewählten Prompt-Engineering-Strategien. Das \textbf{Role-Play Prompting} erfolgt durch die Anweisung \emph{,,Du bist ein erfahrener Barrierefreiheitsexperte für React/TypeScript-Webanwendungen``}, die das Modell in eine domänenspezifische Expertenrolle versetzt (vgl. Abschnitt~\ref{Prompt}). Das \textbf{Few-Shot Prompting} wird durch zwei konkrete Beispiele realisiert: ein \texttt{critical}-Finding (Farbkontrast) mit vollständigem Code-Fix als Must-have-Eintrag und ein \texttt{moderate}-Finding (\texttt{tabIndex}) als Nice-to-have-Eintrag. Beide Beispiele verdeutlichen dem Modell das erwartete Ausgabeformat einschließlich der \ac{WCAG}-Angabe, des Dateipfads und eines konkreten Code-Fix-Vorschlags.

Ergänzend enthält der System-Prompt sechs Priorisierungsregeln, mit denen das Must-have/Nice-to-have-Schema umgesetzt wird:

\begin{enumerate}
  \item Must-have: Severity \texttt{critical}/\texttt{serious} und \ac{WCAG} Level A/AA
  \item Nice-to-have: Severity \texttt{moderate}/\texttt{minor} oder Level AAA
  \item Findings zum selben Element zusammenfassen
  \item Codekontext für konkrete React/TypeScript-Fixes nutzen
  \item Ausschließlich auf Deutsch antworten
  \item Keine Findings erfinden, die nicht in den Daten vorhanden sind
\end{enumerate}

Regel~6 adressiert gezielt das in Abschnitt~\ref{kap:2.4} beschriebene Halluzinationsproblem von \ac{LLM}s. Der vollständige System-Prompt sowie die Few-Shot-Beispiele sind in Anhang~\ref{app:systemprompt} dokumentiert.

\subsection{Token-Budget-Steuerung}

Wie in Abschnitt~\ref{sec:kontextualisierung} beschrieben, ist das Kontextfenster lokaler Modelle eine begrenzte Ressource. Der Prompt-Builder reserviert feste Anteile für den System-Prompt und die Modellausgabe; der verbleibende Platz steht für die serialisierten Findings zur Verfügung.\footnote{Beim eingesetzten Qwen2.5-Coder beträgt das Kontextfenster 32.768~Tokens. Nach Abzug von System-Prompt und Output-Reserve verbleiben ca.~24.000~Tokens für die Findings. Vgl. \cite{hui2024qwen25coder} sowie Anhang \ref{app:tokenbudget}}

Übersteigt die Gesamtmenge der Findings dieses Budget, werden alle Findings quellenübergreifend nach Schweregrad sortiert; Einträge niedriger Priorität werden bei Bedarf entfernt. Findings mit Severity \texttt{critical} und \texttt{serious} haben dabei Vorrang. Der User-Prompt enthält einen expliziten Hinweis, wenn Findings aus Budgetgründen weggelassen werden.

\subsection{Serialisierung}

Jedes Finding wird in ein kompaktes Textformat serialisiert, das die wesentlichen Informationen auf wenige Zeilen verdichtet. Beschreibungen werden auf 300~Zeichen und Code-Snippets auf 20~Zeilen begrenzt, um den Token-Verbrauch je Finding vorhersagbar zu halten. Das vollständige Serialisierungsformat ist im System-Prompt-Beispiel in 
Anhang~\ref{app:systemprompt} illustriert.

% ─────────────────────────────────────────────────────────────────────────────
\section{Ollama-Client und Output-Formatter}
\label{sec:impl-ollama}
\label{sec:impl-output}

Der Ollama-Client kapselt den \acs{HTTP}-Aufruf an die lokale Ollama-Instanz (\texttt{localhost:11434}). Die Anfrage nutzt den \texttt{/api/generate}-Endpunkt mit \texttt{temperature:~0.05} (\ac{NFA}-04).%
\footnote{Der Standardwert \texttt{num\_predict}~=~6\,144 gilt für den operativen Betrieb. In den Benchmarkläufen wurden suiteabhängige Werte eingesetzt: 6\,144 (test-12), 8\,192 (test-50) und 12\,288 (test-100). Vgl. Tabelle~\ref{tab:benchmark_runs} und Abschnitt~\ref{sec:limit_tokenlimit}.}
Die niedrige \texttt{temperature}-Einstellung reduziert den Nicht-Determinismus der Ausgaben (vgl. Abschnitt~\ref{kap:2.4}), ohne ihn vollständig zu eliminieren. Der Client verwendet ein konfigurierbares Timeout (Default: drei Stunden), da die Inferenz auf \ac{CPU}-basierten Systemen je nach Modellgröße mehrere Minuten in Anspruch nehmen kann. Die Antwort von Ollama enthält neben dem generierten Text auch die tatsächliche Anzahl der Prompt- und Output-Tokens (\texttt{prompt\_eval\_count}, \texttt{eval\_count}), die für die Evaluation in Kapitel~\ref{kap:7} erfasst werden.

Der Output-Formatter erzeugt pro Analysedurchlauf zwei Ausgabedateien: einen Markdown-Report für die menschliche Lesbarkeit sowie einen \acs{JSON}-Report für die maschinelle Weiterverarbeitung (\ac{FA}-05). Der Markdown-Report enthält einen Metadaten-Header (\acs{URL}, Modell, Zeitstempel, Token-Statistiken), die Roh-Ausgabe des Sprachmodells sowie eine Metriken-Tabelle mit den Laufzeiten aller Pipeline-Phasen. Der \acs{JSON}-Report speichert zusätzlich das Parsing-Ergebnis der \ac{LLM}-Antwort: Der Formatter versucht, die Markdown-Ausgabe in Must-have- und Nice-to-have-Einträge zu zerlegen. Schlägt dieses Parsing fehl, etwa weil das Modell das vorgegebene Format nicht einhält, wird dies im Report dokumentiert (\texttt{parsed.success: false}).

% ─────────────────────────────────────────────────────────────────────────────
\section{Pipeline-Orchestrierung}
\label{sec:impl-pipeline}

Die Pipeline-Orchestrierung in \texttt{index.ts} verkettet alle Module in der im Anhang dargestellten Reihenfolge (vgl. Anhang~\ref{app:ace-pipeline-schema}, Abbildung~\ref{fig:ace-pipeline-schema}). Jede Phase wird einzeln zeitgemessen, um die in \ac{NFA}-02 geforderte Laufzeitanalyse zu ermöglichen (für die Laufzeitanalyse siehe Kapitel~\ref{kap:7}). Die Metriken werden in einem \texttt{PipelineMetrics}-Objekt gesammelt, das neben den Phasenzeiten auch Token-Schätzungen, Anreicherungsquoten und Deduplizierungsstatistiken enthält.

Der \texttt{-{}-skip-llm}-Modus speichert den fertigen Prompt als Textdatei, ohne Ollama aufzurufen. Dieser Modus wird während der Entwicklung genutzt, um den Prompt unabhängig von der Modellverfügbarkeit zu iterieren und für die Evaluation zu dokumentieren.

\section{Beispielabläufe am Untersuchungsobjekt}
\label{sec:beispielablaeufe}

Der folgende Ablauf zeigt exemplarisch die test-12-Suite und orientiert sich an den acht Pipeline-Phasen aus Abschnitt~\ref{sec:zielarchitektur}:

\begin{description}
  \item[Phase 1 – axe-core] meldet vier Violations, darunter fehlende Alternativtexte und Farbkontrastverstöße.
  \item[Phase 2 – Playwright] Zwei der sieben Interaktionschecks schlagen fehl: \texttt{focus-indicator} (Fokusindikator-Verlust bei Tab-Navigation) und \texttt{focus-after-dialog} (fehlende Fokusübernahme beim Öffnen eines modalen Dialogs).
  \item[Phase 3 – Code-Anreicherung] Grep lokalisiert die zugehörigen Komponenten im Quellcode und ergänzt die relevanten Code-Ausschnitte -- u.\,a. wird \texttt{Modal.tsx} als Ursprungsdatei des \texttt{focus-after-dialog}-Befunds identifiziert.
  \item[Phase 4 – Code-Patterns] Grep erkennt sechs Anti-Pattern-Treffer, darunter \texttt{<div>}-Elemente mit \texttt{onClick}-Handler ohne Tastaturäquivalent.
  \item[Phase 5 – \ac{LLM}-Codeanalyse] Der \ac{LLM}-Detektor analysiert die Quellcode-Chunks und ergänzt je nach Modellgröße durchschnittlich 11 bis 23 weitere semantische Findings.
  \item[Phase 6 – Prompt-Builder] Alle Findings werden serialisiert und das Token-Budget angewendet. Bei \texttt{qwen3:32b} ist das Budget von 6\,144~Tokens ausreichend; bei 7b und 14b ist es gelegentlich knapp.
  \item[Phase 7 – Ollama] Das lokale Modell empfängt den kombinierten Priorisierungsprompt. \texttt{qwen3:32b} klassifiziert den \texttt{focus-after-dialog}-Befund als Must-have und schlägt einen React-konformen \texttt{useEffect}-Hook zur Fokussteuerung vor -- mit Verweis auf die betroffene Datei und Zeilennummer.
  \item[Phase 8 – Output] Der Markdown-Report enthält priorisierte Must-have- und Nice-to-have-Empfehlungen mit Datei- und \ac{WCAG}-Bezug; der \acs{JSON}-Report speichert zusätzlich Token-Statistiken und Phasenlaufzeiten.
\end{description}

Die vollständige Aufstellung der erkannten Violations findet sich in Anhang~\ref{app:accessibility_verstoesse_12}, Screenshots in Anhang~\ref{app:screenshots_test12}, Rohdaten in Anhang~\ref{app:rohdaten_12}.

Auf der test-50-Suite (50 Violations) identifiziert axe-core 14 Violations, zwei Playwright-Checks schlagen fehl und grep liefert 21 Anti-Pattern-Treffer. Die größere Komponentenanzahl erhöht die Token-Last, weshalb das Kontextfenster des 7b-Modells gelegentlich durch Truncation begrenzt wird. Auf der test-100-Suite (100 Violations, 11 Komponenten) verstärkt sich dieser Effekt: das 7b-Modell zeigt regelmäßige Truncation, während \texttt{qwen3:32b} eine stabile Erkennungsleistung beibehält. Screenshots, Rohdaten und Modellvergleich finden sich in den jeweiligen Anhängen (\ref{app:screenshots_test50}, \ref{app:rohdaten_50}, \ref{app:modell_summary_50} bzw.\ \ref{app:screenshots_test100}, \ref{app:rohdaten_100}, \ref{app:modell_summary_100}).

Für den \ac{BWEC} erzeugt die Pipeline zehn axe-core-Violations, zwei Playwright-Checks und acht grep-Treffer. Da die Codebasis weit größer ist, wird sie in Chunks aufgeteilt; bei \texttt{maxfiles}~=~100 (66~Chunks) erzeugt der \acs{LLM}-Detektor im Mittel 15 bis 32 Findings pro Durchlauf, bei \texttt{maxfiles}~=~500 (330~Chunks) entsprechend mehr. Rohdaten und Modellübersichten beider Konfigurationen befinden sich in Anhang~\ref{app:rohdaten_bwec}, \ref{app:modell_summary_bwec} sowie \ref{app:rohdaten_bwec500}, \ref{app:modell_summary_bwec500}. Die vollständige quantitative Auswertung ist in Kapitel~\ref{kap:7} dokumentiert.

Damit bildet der \ac{PoC} die in Kapitel~\ref{kap:5} beschriebene Zielarchitektur vollständig ab: Alle vier Analysequellen sind implementiert, das \ac{UFS} normalisiert die heterogenen Ausgaben, und die zweistufige Prompt-Strategie überführt die kombinierten Findings in entwicklernahe Handlungsempfehlungen. Die Evaluation dieser Implementierung ist Gegenstand des folgenden Kapitels.